\section{Introduction}\label{sec:intro}

Graph neural networks (GNNs)~\citep{kipf2016semi,velickovic2017graph} are used for credit scoring, bail decisions and social recommendation, where the data are split across institutions and the predictions affect people. Federated learning~\citep{mcmahan2017communication} lets such institutions train one model without pooling their graphs. Three requirements then arise together: the model should not discriminate between demographic groups~\citep{dai2021say,agarwal2021towards}; the federation should learn as little as possible about any person's sensitive attribute~\citep{abadi2016deep,mcmahan2018learning}; and a few compromised clients should not be able to steer the outcome~\citep{blanchard2017krum,yin2018byzantine,baruch2019little}.

Group fairness of the global model is a property of the pooled data, so a federation that wants to enforce it needs some statistic of group behaviour from the clients. That statistic is the weak point. A client that reports it can lie, and fair federated methods that reweight or calibrate using client reports~\citep{ezzeldin2023fairfed,li2020fair,fedfact2025} are exposed to such lies by design. In our earlier protocol, \clamp{}, the server saw only four range-checked sums of per-group prediction mass and node counts, computed under verifiable secure aggregation~\citep{bonawitz2017practical,corrigangibbs2017prio}. Reviewers of that work raised four objections, all of which we found to be correct. (1)~One lying client could disable fairness control: a single falsifying client out of ten raised the test gap on Credit from 0.011 to 0.065, about the level of training without control. (2)~The theory was weak: the bound on falsified statistics was vacuous in 10 of 12 attacked settings, the fairness guarantee for the returned model held only on validation data, and one proposition was an accounting identity. (3)~The privacy statement was narrow, and the cap that set the sensitivity was computed from attribute-dependent group sizes in the main experiments while the theorem required a public cap. (4)~Secure aggregation was simulated without any cryptographic cost, and the Gaussian noise of the analysis is not the discrete mechanism a deployment over a finite ring would use.

The first objection cannot be fixed inside the sums-only design. A sum of unverifiable but range-valid inputs moves with each input in proportion to that input's weight, and no range check can tell a plausible lie from the truth. Robustness to a minority of liars requires an order statistic, and an order statistic cannot be computed from sums. This paper therefore changes the trust model of the statistics channel. Each client reports one scalar, its own soft demographic-parity gap on its validation nodes, as additive secret shares to two non-colluding servers, the setting of Prio~\citep{corrigangibbs2017prio} and ELSA~\citep{rathee2023elsa}. The servers check ranges, sort the $K$ gaps, discard the $b$ smallest and the $b$ largest, and release the mean of the rest with Gaussian noise. This released scalar drives the same single-multiplier controller as before. We call the protocol \rt{} (robust, trimmed).

None of the ingredients is new on its own: two-server secret sharing with validity checks is the Prio model, trimmed means are classical robust estimators used for Byzantine-robust learning~\citep{yin2018byzantine}, the Gaussian mechanism with R\'enyi accounting is standard~\citep{mironov2017renyi}, and the multiplier update is projected dual ascent~\citep{cotter2019optimization}. What we contribute is a specific combination for a specific gap: a one-dimensional, privatised and falsification-robust statistics channel for fairness control in federated GNN training, with guarantees that are non-vacuous in the regime we test and with a measured two-party protocol. Concretely:
\begin{enumerate}
\item \textbf{A falsification-robust statistics channel} (Sections~\ref{sec:method} and~\ref{sec:guarantees}; reviewer objections 1 and 2). We prove that with $f\le b$ liars the trimmed mean lies between two order statistics of the honest clients' gaps (Theorem~\ref{thm:falsify}). Liars therefore cannot move the estimate outside the honest range and cannot switch the controller off while most honest clients see a gap above the budget. In experiments with four attack strategies, $f\le b$ liars changed the mean test gap by at most 0.009; the old sums channel lost control with one liar.
\item \textbf{An attribute-level privacy guarantee calibrated by a public floor} (objection 3). Each client's group rates are normalised by $\max(|V_{k,a}|,n_{0,k})$ with $n_{0,k}=\lceil c|V_k|\rceil$, which depends only on the size of the client's validation set. This fixes the sensitivity of the trimmed mean without any attribute-dependent cap, and the experiments use exactly this constant (Theorem~\ref{thm:dp}).
\item \textbf{A measured two-party protocol} (objection 4). We implement the server step in CrypTen~\citep{knott2021crypten} and report time, traffic and rounds; the gap channel costs 0.3--1.1\,s and at most 107\,KB per release for up to 20 clients (Section~\ref{sec:mpc}). The training experiments still run the protocol arithmetic in plaintext with continuous Gaussian noise, and we state what a ring implementation with a discrete mechanism would change.
\item \textbf{An honest account of the price and the limits} (Sections~\ref{sec:results} and~\ref{sec:discussion}). Noise on one trimmed scalar is larger than noise on pooled sums, and with ten releases the honest gap of \rt{} is higher than that of \clamp{} on Credit and the Pokec graphs (0.037--0.061 against 0.018--0.028). On the three small graphs the noise exceeds any plausible gap, so the controller there acts on noise. Robustness ends at the breakdown point $f=b+1$. The trimmed channel does nothing against fairness poisoning through model updates, which the earlier protocol did limit at $P=1$.
\end{enumerate}
We also drop one result of the earlier paper: the a-posteriori record of the dual update was an accounting identity and carried no information about the returned model.

The paper is organised as follows. Section~\ref{sec:related} positions the work. Section~\ref{sec:problem} states the setting and threat model, Section~\ref{sec:method} the protocol and Section~\ref{sec:guarantees} the guarantees, with full proofs in the supplementary material. Sections~\ref{sec:setup} and~\ref{sec:results} report the experiments, and Section~\ref{sec:discussion} the limitations.
